\section{Introduction}\label{sec:introduction}

Let $P=(p_0,\ldots,p_{N-1})$ be a labelled set of points in the plane.
The complete geometric graph on $P$ contains the straight segment $p_ip_j$
for every $i<j$. We assign one of two colours to every edge and count pairs
of edges that cross properly and have the same colour. Both the point set
and the drawing are fixed; the edge colours are the only variables.

The complexity of this problem for complete geometric graphs was posed at
the Crossing Number Workshop in 2019 and appears as Vogtenhuber's Problem~2
in Section~4.12 of the report of Dagstuhl Seminar~21293
\cite{Dagstuhl2021}. Completeness is a substantial restriction. A reduction
that represents a hard graph by segments does not automatically survive
adding every missing segment between their endpoints: the new edges are
free variables and their optimal contribution can depend on the source
colours. Appendix~\ref{app:completion} gives a six-point example of this
dependence. At the other extreme, for points in convex position the optimum
is the known two-page crossing number of a complete graph
\cite{Abrego2013}.

\begin{definition}\label{def:problem}
The input to \emph{complete geometric two-colour crossing minimization}
is a list of pairwise distinct points $P\in(\Q^2)^N$ with no three
collinear, together with an integer $k$ in the decision version.
Coordinates and $k$ are encoded in binary. Write
\[
 E(P)=\{ij:0\leq i<j<N\},\qquad
 M(P,\chi)=\sum_{\{e,f\}\subseteq E(P)}
       \1[e\text{ and }f\text{ cross properly}]\,\1[\chi(e)=\chi(f)]
\]
for $\chi:E(P)\to\{0,1\}$, and put
$\OPT(P)=\min_\chi M(P,\chi)$.
The decision question is whether $\OPT(P)\leq k$.
\end{definition}

A proper crossing is an intersection of two open segments with four
distinct endpoints. Several crossing pairs may have the same geometric
intersection point; each pair is counted separately. The colouring is
unrestricted: it is not a proper edge colouring in the graph-theoretic
sense. There are no weights or prescribed colours.

\begin{theorem}\label{thm:main}
The decision problem in Definition~\ref{def:problem} is NP-complete under
deterministic polynomial-time many-one reductions. Consequently, computing
$\OPT(P)$ exactly is NP-hard. The decision problem remains NP-complete
when all coordinates are integers encoded in binary.
\end{theorem}

We reduce from \textsc{MaxCut} on permutation graphs, with a permutation
representation supplied as part of the input. Its NP-completeness follows
from de Figueiredo et al.\ \cite{Figueiredo2023}: their construction
explicitly produces the two permutations. After relabelling the first
order, a source graph can be described by a permutation $\pi$ of $[n]$,
with an edge $ij$ when $i<j$ and $\pi(i)>\pi(j)$.
For a source graph $G$, set
\[
 g(G)=\min_{x:V(G)\to\{0,1\}}
          \sum_{uv\in E(G)}\1[x(u)=x(v)]
      =|E(G)|-\MaxCut(G).
\]
Thus a source threshold $k$ becomes the equivalent integer threshold
$g(G)\leq t$, where $t=|E(G)|-k$.

Our construction has three interacting components. A large rectangular
comparison grid forces one of two stable colour phases. Small families of
geometric program points then produce controlled fields on the edges
encoding the source. Finally, an exact convex separation removes the
potentially dominant interaction between the grid and program nonfans.
The remaining added edges are included in an explicit error bound. The
conclusion is a uniform interval
\begin{equation}\label{eq:main-sandwich}
 K_{\mathrm{base}}+\Delta g(G)-E_-
 \ \leq\ \OPT(P)\ \leq\
 K_{\mathrm{base}}+\Delta g(G)+E_+,
 \qquad 2^{18}(E_-+E_+)<\Delta,
\end{equation}
where $K_{\mathrm{base}}$ is computable without an optimization oracle.
The lower bound holds for an unrestricted optimum of the complete graph;
the upper bound is attained within the stated error by a single complete
colouring. No colour restrictions are part of the output instance.

\subsection{Conventions and scales}

We use $N$ for the number of points of the final geometric instance.
The letter $N_g$ will denote the number of vertices of a comparison grid
when a distinction is needed. For nonnegative integers $v$, define
\[
 Z(v)=\frac14\left\lfloor\frac v2\right\rfloor
                  \left\lfloor\frac{v-1}{2}\right\rfloor
                  \left\lfloor\frac{v-2}{2}\right\rfloor
                  \left\lfloor\frac{v-3}{2}\right\rfloor
 \quad(v\geq3),\qquad Z(0)=Z(1)=Z(2)=0.
\]
The theorem of \'{A}brego et al.\ \cite[Theorem~7]{Abrego2013} says that
$Z(v)$ is the optimum for a convex complete geometric graph: a same-page
pair of edges in a two-page drawing crosses exactly when its endpoints
alternate in the linear order. Boundary edges do not affect this count.
We use this established value for two convex background components.

For a source with $n\geq1$, the parameters are
\begin{equation}\label{eq:parameters}
\begin{gathered}
 q=2^{20}(n+1),\quad S=nq,\quad d=S+1,\quad H=2^{20000}d^{40},\\
 W=H^4,\quad B=H^5,\quad D_m=512SW,\quad o=2SW+D_m=514SW,\\
 T=2^{13}o^2,\qquad A_0=2^{13}dT,\qquad
 \Delta=q^2W^4.
\end{gathered}
\end{equation}
All exponents and constants are absolute. These deliberately generous
scales simplify exact error separation; no practical implementation of the
full, extremely large target instances is claimed.

\subsection{Membership in NP and integer coordinates}

\begin{proposition}\label{prop:np}
The problem in Definition~\ref{def:problem} belongs to NP. Rational and
integer coordinate versions are polynomial-time equivalent.
\end{proposition}
\begin{proof}
Let $L$ be the total input bit length. Normalize denominators to be
positive and let $D$ be the product of all coordinate denominators.
Its bit length is $O(L)$. Multiplying every coordinate by this same
positive integer yields integer coordinates of $O(L)$ bits and preserves
all orientation signs. Computing $D$ and the scaled coordinates takes
polynomial time by ordinary integer arithmetic.

Pairwise distinctness and absence of collinear triples can be checked
by examining point pairs and the determinants
\[
 \orient(a,b,c)=(b_x-a_x)(c_y-a_y)-(b_y-a_y)(c_x-a_x).
\]
Declare invalid inputs to be no-instances. A certificate consists of the
$\binom N2$ edge colours. For every unordered pair $ab,cd$ with four
distinct endpoints, a proper crossing is equivalent to
\[
 \orient(a,b,c)\orient(a,b,d)<0,
 \qquad \orient(c,d,a)\orient(c,d,b)<0.
\]
There are $O(N^4)$ such tests, each using a constant number of operations
on integers of polynomial bit length. Count those crossing pairs with
equal colours and compare the result with $k$. Negative thresholds are
immediate no-instances. This verifies a certificate in polynomial time.
The same positive scaling is a reduction from rational to integer inputs;
the converse inclusion is immediate.
\end{proof}

The rest of the paper establishes the reduction and the quantitative
interval~\eqref{eq:main-sandwich}. The formalization and its precise
external dependencies are described in Section~\ref{sec:formalization}.
